\chapter{Encuesta aplicada a docentes de la asignatura FLP}
\label{annex:encuesta-docentes}

El siguiente instrumento fue aplicado a 3 docentes con experiencia en la impartición de la asignatura \textit{Fundamentos de Interpretación y Compilación de Lenguajes de Programación} en la Universidad del Valle, sede Tuluá. Los datos consignados corresponden a estimaciones obtenidas a partir del contacto con el cuerpo docente durante el desarrollo del proyecto.

\section*{Instrucciones}

Por favor responda con base en su experiencia docente en la asignatura FLP. Las respuestas son confidenciales y tienen como fin identificar las principales dificultades de los estudiantes para orientar el diseño de una herramienta educativa de apoyo.

\section*{Bloque I — Perfil docente}

\begin{enumerate}
    \item Número de semestres que ha dictado la asignatura FLP:
    \begin{enumerate}
        \item 1 – 2 semestres
        \item 3 – 5 semestres
        \item Más de 5 semestres
    \end{enumerate}

    \item Enfoque metodológico principal que utiliza en la asignatura:
    \begin{enumerate}
        \item Clases magistrales con ejercicios guiados.
        \item Proyectos de implementación progresiva (intérpretes/compiladores).
        \item Combinación de ambas.
        \item Otro: \underline{\hspace{6cm}}
    \end{enumerate}
\end{enumerate}

\section*{Bloque II — Dificultades observadas en los estudiantes}

\begin{enumerate}
    \setcounter{enumi}{2}

    \item ¿Ha observado deficiencias generales en los estudiantes que ingresan a la asignatura?
    \begin{enumerate}
        \item Sí, de manera consistente.
        \item En ocasiones.
        \item No.
    \end{enumerate}

    \item ¿En cuáles cursos previos observa mayores falencias que afectan el desempeño en FLP? \textit{(puede marcar varias)}:
    \begin{enumerate}
        \item Matemáticas discretas (lógica proposicional, relaciones, conjuntos).
        \item Teoría de lenguajes y autómatas.
        \item Estructuras de datos (árboles, pilas, recursión).
        \item Programación funcional.
        \item Otro: \underline{\hspace{6cm}}
    \end{enumerate}

    \item Indique el nivel de dificultad que observa en sus estudiantes para cada concepto de la asignatura, usando la escala:

    \begin{center}
    \textbf{1 = Ninguna \quad 2 = Leve \quad 3 = Moderada \quad 4 = Alta \quad 5 = Muy alta}
    \end{center}

    \begin{longtable}{|p{9cm}|p{0.5cm}|p{0.5cm}|p{0.5cm}|p{0.5cm}|p{0.5cm}|}
        \hline
        \textbf{Concepto} & \textbf{1} & \textbf{2} & \textbf{3} & \textbf{4} & \textbf{5} \\
        \hline
        Definición y lectura de gramáticas libres de contexto & & & & & \\ \hline
        Análisis sintáctico (\textit{parsing}) y construcción de derivaciones & & & & & \\ \hline
        Construcción e interpretación de árboles de sintaxis abstracta (AST) & & & & & \\ \hline
        Gestión de ambientes de ejecución y alcance léxico & & & & & \\ \hline
        Clausuras y funciones de orden superior & & & & & \\ \hline
        Definiciones dirigidas por sintaxis (SDD) & & & & & \\ \hline
        \caption{Instrumento — Dificultad observada por docentes en estudiantes de FLP.\\
        \scriptsize \textbf{Fuente:} Elaboración propia.}
        \label{tab:instrumento-docentes}
    \end{longtable}
\end{enumerate}

\section*{Bloque III — Herramientas y estrategias didácticas}

\begin{enumerate}
    \setcounter{enumi}{5}

    \item ¿Ha utilizado estrategias o materiales adicionales para apoyar la comprensión de los conceptos más abstractos de la asignatura?
    \begin{enumerate}
        \item Sí, con buenos resultados.
        \item Sí, pero con resultados limitados.
        \item No, no he encontrado recursos adecuados.
    \end{enumerate}

    \item En su opinión, ¿cuál es la principal limitación de Racket y la librería EOPL como herramientas pedagógicas?
    \begin{enumerate}
        \item La curva de aprendizaje técnica consume tiempo que podría usarse en los conceptos.
        \item No exponen visualmente los procesos internos de interpretación.
        \item Los estudiantes sin experiencia previa en programación funcional tienen dificultades para configurarlas.
        \item Otra: \underline{\hspace{6cm}}
    \end{enumerate}
\end{enumerate}

\section*{Bloque IV — Valoración de herramientas interactivas}

\begin{enumerate}
    \setcounter{enumi}{7}

    \item ¿Considera que una herramienta web interactiva que permita definir gramáticas, visualizar el AST y representar los cambios en el ambiente de ejecución mejoraría la comprensión de los estudiantes?
    \begin{enumerate}
        \item Sí, significativamente.
        \item Probablemente sí.
        \item No lo sé.
        \item Probablemente no.
    \end{enumerate}

    \item ¿Qué características considera imprescindibles en dicha herramienta? \textit{(puede marcar varias)}:
    \begin{enumerate}
        \item Acceso desde el navegador sin instalación.
        \item Ejemplos predefinidos de gramáticas y programas de menor a mayor complejidad.
        \item Visualización gráfica del AST con posibilidad de explorar cada nodo.
        \item Representación paso a paso de los cambios en el ambiente de ejecución.
        \item Retroalimentación inmediata ante errores sintácticos o semánticos.
        \item Otra: \underline{\hspace{6cm}}
    \end{enumerate}

    \item Comentarios adicionales sobre las dificultades que observa y los recursos que harían falta: \underline{\hspace{12cm}}

    \underline{\hspace{14cm}}
\end{enumerate}
